% Generated draft: requires mathematical and source review.
\documentclass[11pt]{article}
\usepackage[margin=25mm]{geometry}
\usepackage{fontspec}
\setmainfont{Latin Modern Roman}
\usepackage{amsmath,amssymb,mathtools,unicode-math}
\setmathfont{Latin Modern Math}
\usepackage{enumitem,needspace,xurl,hyperref}
\setlength{\parindent}{0pt}
\setlength{\parskip}{5pt}
\setlength{\emergencystretch}{3em}
\newcommand{\N}{\mathbb{N}}
\newcommand{\Z}{\mathbb{Z}}
\newcommand{\Q}{\mathbb{Q}}
\newcommand{\R}{\mathbb{R}}
\newcommand{\C}{\mathbb{C}}
\newcommand{\F}{\mathbb{F}}
\newcommand{\A}{\mathbb{A}}
\newcommand{\PP}{\mathbb{P}}
\newcommand{\E}{\mathbb{E}}
\newcommand{\eps}{\varepsilon}
\newcommand{\dd}{\,\mathrm{d}}
\begin{document}
\section*{UnsolvedMath ID: 800829}
\textbf{Source problem code:} catalogue entry
\subsection{Definitions and mathematical statement}
Let \(M\) denote the three-dimensional flow domain and let \(v\) denote a steady velocity field on \(M\). Write \(\omega=\operatorname{curl}v\) for its vorticity. The geometry defining curl, the governing steady-flow equation, the admissible regularity class, and any boundary conditions are not specified in the supplied record; these must be fixed before a precise genericity statement can be made.

The source distinguishes fields that are not everywhere collinear with \(\omega\), fields satisfying \(\omega=fv\) for a scalar function \(f\), and curl eigenfields satisfying
\[
\operatorname{curl}v=\lambda v,
\qquad \lambda\text{ constant},\qquad v\not\equiv0.
\]
It calls the last class Beltrami fields. The representation \(\omega=fv\) requires care at zeros of \(v\); the source supplies no nonvanishing assumption or regularity requirement on \(f\).

An invariant torus or annulus means a surface of the corresponding topological type that is preserved by the flow of \(v\). The source attributes to Arnold the assertion that an analytic steady field not everywhere collinear with its curl has a flow domain almost everywhere fibered into invariant tori and annuli. It further asserts a similar structure when \(\operatorname{curl}v=fv\) and \(f\) is a generic function, whereas constant-coefficient curl eigenfields can have intricate topology. These are supplied contextual assertions: the precise hypotheses, the meaning of genericity for \(f\), and the measure or exceptional set intended by ``almost everywhere'' are not provided.

\textbf{Problem.} Explain the apparent tension between the regular topology attributed to generic steady fields and the intricate topology possible for curl eigenfields. What should constitute a typical steady field, and what notion of genericity is natural on the space of steady fields?

A precise formulation must identify the admissible space \(\mathcal S\) of steady fields, its geometric and boundary assumptions, and the topology, measure, or other structure used to define typicality. It must also specify the relevant meaning of regular or chaotic flow topology. The supplied source does not select these data. In particular, its statement that a sufficiently chaotic steady field must be a curl eigenfield is part of the motivating discussion, not a fully quantified implication with supplied hypotheses and a definition of sufficient chaos.
\subsection{Short English statement}
The problem asks how to define a typical steady three-dimensional fluid flow so as to explain why the source describes generic steady flows as topologically regular, while curl eigenfields can have complicated dynamics. The flow domain, admissible steady equations, boundary conditions, and meaning of genericity remain unspecified. This is a statement-only draft requiring mathematical review; the registry's open status is supplied metadata.
\subsection{Sources}
\begin{itemize}
\item Supplied local UnsolvedMath registry record: permanent ID \(800829\), exact source problem code \texttt{catalogue entry}, title ``Geometric Hydrodynamics in Open Problems --- Problem 28.'' The record attributes the problem to Khesin, Misiolek, and Modin, with proposed year 2022, and labels it open. This status has not been independently verified.
\item Khesin, Misiolek, and Modin, \emph{Geometric Hydrodynamics in Open Problems} (2022), Problem 28. Supplied source URL: \url{https://arxiv.org/abs/2205.01143}. The background reports access on 2026-10-16 and extraction by \texttt{arxiv\_latex}; these are supplied provenance details, not a claim of independent source consultation.
\item The supplied statement attributes the invariant-surface discussion to Arnold and cites the reference labels \texttt{[arn66]} and \texttt{[ak]}. Their bibliographic entries are not included in the supplied record, so their identities and precise theorem hypotheses remain unresolved.
\end{itemize}
\end{document}
